% Part: first-order-logic
% Chapter: sequent-calculus
% Section: proving-things

\documentclass[../../../include/open-logic-section]{subfiles}

\begin{document}

\iftag{FOL}
      {\olfileid{fol}{seq}{pro}}
      {\olfileid{pl}{seq}{pro}}

\olsection{د \usetoken{P}{derivation} بېلګې}

\begin{ex}
د سېکوېنټ $!A \land !B \Sequent !A$ يو
$\Log{LK}$-!!{derivation} ورکړئ۔

د !!{derivation} په لاندې برخه کښې د غوښتل شوي وروستي سېکوېنټ په
ليکلو پيل کوو۔
\begin{prooftree}
\AxiomC{}
\UnaryInf$!A\land !B \fCenter !A$
\end{prooftree}
بيا بايد معلومه کړو چې د کوم ډول استنتاج لاندې سېکوېنټ دا بڼه
لرلے شي۔ دا يوه جوړښتي قاعده کېدے شي، خو غوره ده چې لومړے کومه
منطقي قاعده ولټوؤ۔ په لاندې سېکوېنټ کښې يوازينے منطقي نښلوونکے
$\land$ دے؛ نو د $\land$ قاعده لټوؤ، او ځکه چې د $\land$ نښه په
مقدم کښې راځي، د \LeftR{\land} قاعده ګورو۔
\begin{prooftree}
\AxiomC{}
\RightLabel{\LeftR{\land}}
\UnaryInf$!A\land !B \fCenter !A$
\end{prooftree}
د \LeftR{\land} استنتاج د پورتني سېکوېنټ دپاره دوه امکانونه شته:
يا $!A \Sequent !A$ او يا $!B \Sequent !A$۔ څرګنده ده چې
$!A \Sequent !A$ يو لومړنے سېکوېنټ دے، او دا ښه خبره ده؛ خو
$!B \Sequent !A$ په عمومي توګه د اشتقاق وړ نۀ دے۔ پورتنے سېکوېنټ
ورډکوو:
\begin{prooftree}
\Axiom$!A \fCenter !A$
\RightLabel{\LeftR{\land}}
\UnaryInf$!A\land !B \fCenter !A$
\end{prooftree}
اوس د سېکوېنټ $!A \land !B \Sequent !A$ سم
$\Log{LK}$-!!{derivation} لرو۔
\end{ex}

\begin{ex}
د سېکوېنټ $\lnot !A \lor !B \Sequent !A \lif !B$ يو
$\Log{LK}$-!!{derivation} ورکړئ۔

د !!{derivation} په لاندې برخه کښې د غوښتل شوي وروستي سېکوېنټ په
ليکلو پيل وکړئ۔
\begin{prooftree}
\AxiomC{}
\UnaryInf$\lnot !A \lor !B \fCenter !A \lif !B$
\end{prooftree}
د داسې منطقي قاعدې د موندلو دپاره چې دا وروستنے سېکوېنټ راکړي، د
وروستي سېکوېنټ منطقي نښلوونکي ګورو: $\lnot$، $\lor$ او $\lif$۔
اوس مهال يوازې $\lor$ او $\lif$ مهم دي، ځکه چې د وروستي سېکوېنټ د
!!{sentence}s !!{main operator}s دي؛ خو $\lnot$ د بل نښلوونکي د
چاپېريال دننه دے، نو وروسته به ئې وڅېړو۔ له دې امله د وروستي
استنتاج د منطقي قاعدې امکانونه \LeftR{\lor} او \RightR{\lif} دي۔
په رښتيا هره يوه غوره کولے شو، خو \RightR{\lif} غوره کوو، يوازې
د دې دپاره چې په دوو څانګو وېشل لږ وروسته کړو۔ د هغو
\RightR{\lif} استنتاجونو د بڼې له مخې چې دا لاندې سېکوېنټ ورکولے
شي، بڼه بايد داسې وي:
\begin{prooftree}
\AxiomC{}
\UnaryInf$ !A, \lnot !A \lor !B \fCenter !B $
\RightLabel{\RightR{\lif}} \UnaryInf$ \lnot !A \lor !B \fCenter !A \lif !B $
\end{prooftree}

که $\lnot !A \lor !B$ د مقدم بهرني پاے ته يوسو، د
\LeftR{\lor} قاعده کارولے شو۔ د شېما له مخې دا بايد په لاندې ډول
پر دوو پورتنيو سېکوېنټونو ووېشل شي۔
\textbf{د ژباړې د نسخې سمون (OLSQ-001):} د سرچينې په هغو څلورو
قاعدوي ونو کښې چې د مقدم دوه فارمولونه سره بدلوي، د ښي لوري د
تبادلې نښه راغلې ده؛ دلته د قاعدې له واقعي لوري سره سمه د کيڼ
لوري د تبادلې نښه کارول کېږي۔
\begin{prooftree}
\AxiomC{}
\UnaryInf$\lnot !A, !A \fCenter !B$
\AxiomC{}
\UnaryInf$!B, !A \fCenter !B$
\RightLabel{\LeftR{\lor}}
\BinaryInf$ \lnot !A \lor !B, !A \fCenter !B $
\RightLabel{\LeftR{\Exchange}}
\UnaryInf$ !A, \lnot !A \lor !B \fCenter !B $
\RightLabel{\RightR{\lif}}
\UnaryInf$ \lnot !A \lor !B \fCenter !A \lif !B $
\end{prooftree}
راياد کړئ چې موږ هڅه کوو تر لومړنيو سېکوېنټونو پورې په پورته لوري
لار پيدا کړو؛ داسې ښکاري چې ډېر ورنژدې يو!{} ښۍ څانګه له لومړني
سېکوېنټ څخه يوازې يو کمزوري کول او يوه تبادله لرې ده، او بيا
بشپړېږي:
\begin{prooftree}
\AxiomC{}
\UnaryInf$\lnot !A, !A \fCenter !B$
\Axiom$!B \fCenter !B$
\RightLabel{\LeftR{\Weakening}}
\UnaryInf$!A, !B \fCenter !B$
\RightLabel{\LeftR{\Exchange}}
\UnaryInf$!B, !A \fCenter !B$
\RightLabel{\LeftR{\lor}}
\BinaryInf$\lnot !A \lor !B, !A \fCenter !B $
\RightLabel{\LeftR{\Exchange}}
\UnaryInf$ !A, \lnot !A \lor !B \fCenter !B $
\RightLabel{\RightR{\lif}}
\UnaryInf$ \lnot !A \lor !B \fCenter !A \lif !B $
\end{prooftree}

اوس کيڼه څانګه ګورو۔ په هره !!{sentence} کښې يوازينے منطقي
نښلوونکے د مقدم د !!{sentence}s د $\lnot$ نښه ده، نو د
\LeftR{\lnot} قاعدې يوه کارونه ګورو۔
\begin{prooftree}
\AxiomC{}
\UnaryInf$ !A \fCenter !B, !A$
\RightLabel{\LeftR{\lnot}}
\UnaryInf$\lnot !A, !A \fCenter !B$
\Axiom$!B \fCenter !B$
\RightLabel{\LeftR{\Weakening}}
\UnaryInf$!A, !B \fCenter !B$
\RightLabel{\LeftR{\Exchange}}
\UnaryInf$!B, !A \fCenter !B$
\RightLabel{\LeftR{\lor}}
\BinaryInf$\lnot !A \lor !B, !A \fCenter !B $
\RightLabel{\LeftR{\Exchange}}
\UnaryInf$ !A, \lnot !A \lor !B \fCenter !B $
\RightLabel{\RightR{\lif}}
\UnaryInf$ \lnot !A \lor !B \fCenter !A \lif !B $
\end{prooftree}
لکه ښۍ څانګه چې مو بشپړه کړه، له دې کيڼې څانګې څخه هم يوازې يو
کمزوري کول او يوه تبادله لرې يو۔
\begin{prooftree}
\Axiom$!A \fCenter !A$
\RightLabel{\RightR{\Weakening}}
\UnaryInf$ !A \fCenter !A, !B$
\RightLabel{\RightR{\Exchange}}
\UnaryInf$ !A \fCenter !B, !A$
\RightLabel{\LeftR{\lnot}}
\UnaryInf$\lnot !A, !A \fCenter !B$
\Axiom$!B \fCenter !B$
\RightLabel{\LeftR{\Weakening}}
\UnaryInf$!A, !B \fCenter !B$
\RightLabel{\LeftR{\Exchange}}
\UnaryInf$!B, !A \fCenter !B$
\RightLabel{\LeftR{\lor}}
\BinaryInf$\lnot !A \lor !B, !A \fCenter !B $
\RightLabel{\LeftR{\Exchange}}
\UnaryInf$ !A, \lnot !A \lor !B \fCenter !B $
\RightLabel{\RightR{\lif}}
\UnaryInf$ \lnot !A \lor !B \fCenter !A \lif !B $
\end{prooftree}
\end{ex}

\begin{ex}
د سېکوېنټ $\lnot !A \lor \lnot !B
\Sequent \lnot (!A \land !B)$ يو $\Log{LK}$-!!{derivation} ورکړئ۔

د پورته طريقو په کارولو، غوښتل شوے وروستنے سېکوېنټ په لاندې برخه
کښې ليکو۔
\begin{prooftree}
\AxiomC{}
\UnaryInf$ \lnot !A \lor \lnot !B \fCenter \lnot (!A \land !B) $
\end{prooftree}
په وروستي سېکوېنټ کښې د !!{sentence}s شته اصلي نښلوونکي د $\lor$
او $\lnot$ نښې دي۔ دلته هم \LeftR{\lor} او هم \RightR{\lnot}
کار ورکوي، خو په \RightR{\lnot} پيل کوو، ځکه چې د يوې شېبې دپاره
پر دوو څانګو وېشل ځنډوي:
\begin{prooftree}
\AxiomC{}
\UnaryInf$!A \land !B, \lnot !A \lor \lnot !B \fCenter $
\RightLabel{\RightR{\lnot}}
\UnaryInf$\lnot !A \lor \lnot !B \fCenter \lnot (!A \land !B)$
\end{prooftree}
اوس د \LeftR{\land} او \LeftR{\lor} تر منځ ټاکنه لرو۔ ګورو چې د
\LeftR{\land} قاعدې په کارولو څه کېږي: يا په سېکوېنټ $!A,
\lnot !A \lor \lnot !B \Sequent \quad$ او يا په سېکوېنټ $!B, \lnot !A
\lor \lnot !B \Sequent \quad$ پيل کولے شو۔ ځکه چې !!{derivation} د
$!A$ او $!B$ له پلوه متناظر دے، لومړنے غوره کوو۔
\textbf{د ژباړې د نسخې سمون (OLSQ-002):} د سرچينې په دغو دوو
امکاني پورتنيو سېکوېنټونو کښې له دويمې جملې څخه نفي پاتې شوې وه؛
دلته د وروستي سېکوېنټ او ورپسې قاعدوي ونې له مخې په دواړو جملو
کښې نفي ساتل کېږي۔
\begin{prooftree}
\AxiomC{}
\UnaryInf$!A, \lnot !A \lor \lnot !B \fCenter $
\RightLabel{\LeftR{\land}}
\UnaryInf$!A \land !B, \lnot !A \lor \lnot !B \fCenter $
\RightLabel{\RightR{\lnot}}
\UnaryInf$\lnot !A \lor \lnot !B \fCenter \lnot (!A \land !B)$
\end{prooftree}
د !!{derivation} د ډکولو په دوام وينو چې له يوې ستونزې سره مخ کېږو:
\begin{prooftree}
\Axiom$!A \fCenter !A$
\RightLabel{\LeftR{\lnot}}
\UnaryInf$ \lnot !A, !A \fCenter$
\AxiomC{}
\RightLabel{?}
\UnaryInf$!A \fCenter !B$
\RightLabel{\LeftR{\lnot}}
\UnaryInf$ \lnot !B, !A \fCenter$
\RightLabel{\LeftR{\lor}}
\BinaryInf$\lnot !A \lor \lnot !B, !A \fCenter $
\RightLabel{\LeftR{\Exchange}}
\UnaryInf$!A, \lnot !A \lor \lnot !B \fCenter $
\RightLabel{\LeftR{\land}}
\UnaryInf$!A \land !B, \lnot !A \lor \lnot !B \fCenter $
\RightLabel{\RightR{\lnot}}
\UnaryInf$\lnot !A \lor \lnot !B \fCenter \lnot (!A \land !B)$
\end{prooftree}
د ښۍ څانګې پورته برخه نوره نۀ شي راکمېدلے او د جوړښتي استنتاجونو
له لارې لومړني سېکوېنټ ته نۀ شي رسول کېدے؛ نو دا سمه لار نۀ ده۔
له دې امله پورته د \LeftR{\land} قاعدې کارول تېروتنه وه۔ بېرته
مخکيني حالت ته ځو او پر ځاے ئې \LeftR{\lor} قاعده پلې کوو، نو
ترلاسه کوو:
\begin{prooftree}
\AxiomC{}
\UnaryInf$\lnot !A, !A \land !B \fCenter $

\AxiomC{}
\UnaryInf$\lnot !B, !A \land !B \fCenter $

\RightLabel{\LeftR{\lor}}
\BinaryInf$\lnot !A \lor \lnot !B, !A \land !B \fCenter $
\RightLabel{\LeftR{\Exchange}}
\UnaryInf$!A \land !B, \lnot !A \lor \lnot !B \fCenter $
\RightLabel{\RightR{\lnot}}
\UnaryInf$\lnot !A \lor \lnot !B \fCenter \lnot (!A \land !B)$
\end{prooftree}
هره څانګه د مخکښې په شان بشپړه کړو، نو ترلاسه کوو:
\begin{prooftree}
\Axiom$ !A \fCenter!A$
\RightLabel{\LeftR{\land}}
\UnaryInf$!A \land !B \fCenter !A$
\RightLabel{\LeftR{\lnot}}
\UnaryInf$\lnot !A, !A \land !B \fCenter $

\Axiom$ !B \fCenter !B$
\RightLabel{\LeftR{\land}}
\UnaryInf$!A \land !B \fCenter !B$
\RightLabel{\LeftR{\lnot}}
\UnaryInf$\lnot !B, !A \land !B \fCenter $

\RightLabel{\LeftR{\lor}}
\BinaryInf$\lnot !A \lor \lnot !B, !A \land !B \fCenter $
\RightLabel{\LeftR{\Exchange}}
\UnaryInf$!A \land !B, \lnot !A \lor \lnot !B \fCenter $
\RightLabel{\RightR{\lnot}}
\UnaryInf$\lnot !A \lor \lnot !B \fCenter \lnot (!A \land !B)$
\end{prooftree}
(په دغو ګامونو کښې د $\land$ قاعدې د $\lnot$ له قاعدو څخه لاندې
هم پلې کېدے شوې او بيا به هم سم !!{derivation} ترلاسه شوے و۔)
\end{ex}

\begin{ex}
تر اوسه مو د انقباض قاعده نۀ ده کارولې، خو کله کله اړينه وي۔ دا
ئې يوه بېلګه ده۔ فرض کړئ غواړو $\quad \Sequent !A \lor \lnot !A$
ثابت کړو۔ د $\RightR{\lor}$ سرچپه پلي کول به له دغو دوو
!!{derivation}s څخه يو راکړي:
\begin{prooftree}
\AxiomC{}
\UnaryInf$ \fCenter !A$
\RightLabel{\RightR{\lor}}
\UnaryInf$ \fCenter !A \lor \lnot !A$
\DisplayProof\qquad\bottomAlignProof
\AxiomC{}
\UnaryInf$!A \fCenter $
\RightLabel{\RightR{\lnot}}
\UnaryInf$ \fCenter \lnot !A$
\RightLabel{\RightR{\lor}}
\UnaryInf$ \fCenter !A \lor \lnot !A$
\end{prooftree}
البته له دغو څخه هېڅ يو په لومړني سېکوېنټ نۀ پاے ته رسېږي۔ چل دا
دے چې پوه شو د انقباض قاعده د !!a{sentence} دوه نقلونه سره يو
کولے شي---او کله چې د ثبوت لټون کوو، يعنې له لاندې څخه پورته ځو،
د $!A \lor \lnot !A$ يوه کاپي په مقدمه کښې ساتلے شو، مثلاً:
\begin{prooftree}
\AxiomC{}
\UnaryInf$ \fCenter !A \lor \lnot !A, !A$
\RightLabel{\RightR{\lor}}
\UnaryInf$ \fCenter !A \lor \lnot !A, !A \lor \lnot !A$
\RightLabel{\RightR{\Contraction}}
\UnaryInf$ \fCenter !A \lor \lnot !A$
\end{prooftree}
اوس $\RightR{\lor}$ دويم ځل کارولے شو او~$\lnot !A$ هم ترلاسه کوو،
چې بشپړ !!{derivation} ته رسېږي۔
\begin{prooftree}
\Axiom$!A \fCenter !A$
\RightLabel{\RightR{\lnot}}
\UnaryInf$\fCenter !A, \lnot !A$
\RightLabel{\RightR{\lor}}
\UnaryInf$\fCenter !A, !A \lor \lnot !A$
\RightLabel{\RightR{\Exchange}}
\UnaryInf$ \fCenter !A \lor \lnot !A, !A$
\RightLabel{\RightR{\lor}}
\UnaryInf$ \fCenter !A \lor \lnot !A, !A \lor \lnot !A$
\RightLabel{\RightR{\Contraction}}
\UnaryInf$ \fCenter !A \lor \lnot !A$
\end{prooftree}
\end{ex}

\begin{prob}
د لاندې سېکوېنټونو !!{derivation}s ورکړئ:
\begin{enumerate}
\item $!A \land (!B \land !C) \Sequent (!A \land !B) \land !C$.
\item $!A \lor (!B \lor !C) \Sequent (!A \lor !B) \lor !C$.
\item $!A \lif (!B \lif !C) \Sequent !B \lif (!A \lif !C)$.
\item $!A \Sequent \lnot\lnot !A$.
\end{enumerate}
\end{prob}

\begin{prob}
د لاندې سېکوېنټونو !!{derivation}s ورکړئ:
\begin{enumerate}
\item $(!A \lor !B) \lif !C \Sequent !A \lif !C$.
\item $(!A \lif !C) \land (!B \lif !C) \Sequent (!A \lor !B) \lif !C$.
\item $\Sequent \lnot(!A \land \lnot !A)$.
\item $!B \lif !A \Sequent \lnot !A \lif \lnot !B$.
\item $\Sequent (!A \lif \lnot !A) \lif \lnot !A$.
\item $\Sequent \lnot(!A \lif !B) \lif \lnot !B$.
\item $!A \lif !C \Sequent \lnot (!A \land \lnot !C)$.
\item $!A \land \lnot !C \Sequent \lnot (!A \lif !C)$.
\item $!A \lor !B, \lnot !B \Sequent !A$.
\item $\lnot !A \lor \lnot !B \Sequent \lnot(!A \land !B)$.
\item $\Sequent (\lnot !A \land \lnot !B) \lif\lnot(!A \lor !B)$.
\item $\Sequent \lnot(!A \lor !B) \lif (\lnot !A \land \lnot !B)$.
\end{enumerate}
\end{prob}

\begin{prob}
د لاندې سېکوېنټونو !!{derivation}s ورکړئ:
\begin{enumerate}
\item $\lnot(!A \lif !B) \Sequent !A$.
\item $\lnot(!A \land !B) \Sequent \lnot !A \lor \lnot !B$.
\item $!A \lif !B \Sequent \lnot !A \lor !B$.
\item $\Sequent \lnot \lnot !A \lif !A$.
\item $!A \lif !B, \lnot !A \lif !B \Sequent !B$.
\item $(!A \land !B) \lif !C \Sequent (!A \lif !C) \lor (!B \lif !C)$.
\item $(!A \lif !B) \lif !A \Sequent !A$.
\item $\Sequent (!A \lif !B) \lor (!B \lif !C)$.
\end{enumerate}
(دا ټول د~$\RightR{\Contraction}$ قاعدې ته اړتيا لري۔)
\end{prob}
\end{document}
